\documentclass[10pt]{article}
\usepackage[french]{babel}
\usepackage[utf8]{inputenc}
\usepackage[T1]{fontenc}
\usepackage{amsmath}
\usepackage{amsfonts}
\usepackage{amssymb}
\usepackage[version=4]{mhchem}
\usepackage{stmaryrd}

\begin{document}
\section*{Analyse Numérique - TD n \({ }^{\circ} \mathbf{1}\) 05-02-2025 - SIGMA1}
\section*{Exercice 1:}
On considère la fonction \(f(x)=\sin (\pi x), x \in[0,1]\). Soient \(P_{2}\) le polynôme d'interpolation de degré 2 de \(f\) aux points \(x_{0}=0, x_{1}=1 / 2, x_{2}=1\) et ( \(P_{21}, P_{22}\) ) les polynômes interpolant \(f\) par intervalles par deux polynômes de degré 2 dans le cas où \([0,1]\) est subdivisé en deux intervalles \([0,1 / 2]\) et \([1 / 2,1]\) de longueur \(h=1 / 2\). Une fois trouvés \(P_{2}\) et \(P_{21}\), par exemple par la méthode de Lagrange, calculer \(f(1 / 8), P_{2}(1 / 8)\) et \(P_{21}(1 / 8)\).

\section*{Exercice 2:}
On considère l'interpolation par polynômes de la fonction \(\sin u s(x), x\) exprimé en radians, aux points \(\left(\frac{\pi}{2}, \frac{2 \pi}{3}, \frac{3 \pi}{4}, \frac{5 \pi}{6}\right)\). La méthode des différences divisées est utilisée.

1- Trouver le polynôme \(P(x)\). Les grandeurs calculées sont représentées par leur valeur décimale ;\\
2- Estimer l'erreur \(|P(x)-\sin (x)|\) sur les points d'interpolation ;\\
3- Estimer l'erreur \(|P(x)-\sin (x)|\) lorsque \(x=\pi\). Que déduisez-vous sur le comportement de l'erreur hors de l'intervalle d'interpolation ?

\section*{Exercice 3:}
On pose

\[
A=\left(\begin{array}{cccc}
-2 & 1 & -1 & 1 \\
2 & 0 & 4 & -3 \\
-4 & -1 & -12 & 9 \\
-2 & 1 & 1 & -4
\end{array}\right)
\]

\begin{enumerate}
  \item Donner la décomposition \(L U\) de la matrice \(A\) (i.e \(A=L U\) avec \(L\) triangulaire inférieure et \(U\) triangulaire supérieure)
  \item En déduire la solution du système linéaire \(A \mathrm{x}=b\) où \(\mathrm{b}=(1.5,4,-14,-6.5)^{\top}\).
\end{enumerate}

\section*{Exercice 4:}
Soit à résoudre le système suivant :\\
\(\left\{\begin{array}{c}4 x_{1}+2 x_{2}+x_{3}=11 \\ -x_{1}+2 x_{2}=3 \\ 2 x_{1}+x_{2}+4 x_{3}=16\end{array}\right.\)\\
On sait que ce système admet pour solution : \(x_{1}=1 ; x_{2}=2 ; x_{3}=3\)

\begin{enumerate}
  \item Ecrire le système sous forme matricielle \(\mathrm{AX}=\mathrm{B}\) et utiliser la Méthode de Jacobi pour le résoudre. Vérifier la condition de diagonale dominante pour la matrice \(A\).
  \item Utiliser à présent La méthode de Gauss-Seidel pour résoudre le système
\end{enumerate}

\end{document}